Three hundred problems where number theory, algebra, probability and analysis meet. Why do Fibonacci numbers repeat
modulo a prime with a period that divides $p-1$ or $2(p+1)$? Which strange pair of dice rolls every sum with the same odds as
an ordinary pair? Why is $\zeta(3)$ irrational, and why must a set of points with no three on a line be so small in a space over
the field with three elements? Each problem mixes ideas from at least two volumes of the first part of \emph{Order and Chance}, and none says which
tool to use: finding it is the game. Hints point the way when the search stalls, and every solution begins with how one
might discover the approach before giving a complete proof. Long \emph{Journeys} walk step by step to famous theorems, from
Lucas--Lehmer to the cap set bound, and remarks show where each problem leads in current research. Companion Python
notebooks check every number and turn the ideas into experiments and algorithms. A book for anyone who has learned the
theorems and now wants to use them.
